% Section: Cold-Start Prior Value (K=2)
% Include from paper/sections/results.tex

\subsection{The Value of Warmup Priors ($K{=}2$)}
\label{sec:warmup_ablation_k2}

\paragraph{Motivation.}
The Pareto frontier in Figure~\ref{fig:pareto_k2} evaluates routing
quality after extensive online learning ($n{=}8{,}374$ training
prompts), a regime where both warm-started and uninformed bandits
converge to similar performance.
In production, however, the highest-stakes period is
\emph{deployment day one}: when the router goes live with zero
online experience, every suboptimal routing decision incurs
real cost and latency.
This section isolates the cold-start value proposition of warmup
priors by evaluating from the very first prompt.

\paragraph{Conditions.}
We compare two conditions, both using the production hyperparameters
from Figures~\ref{fig:pareto_k2} and~\ref{fig:adaptive_drift}:

\begin{itemize}[leftmargin=*,noitemsep]
  \item \textbf{BanditGPT (warm priors):}  Disjoint LinUCB initialised
        with offline priors ($n_{\mathrm{eff}}{=}5{,}000$,
        $\alpha{=}0.5$).
  \item \textbf{Tabula Rasa (no priors):}  Same architecture and
        exploration rate ($\alpha{=}0.5$), starting from scratch
        ($A = \lambda I$, $b = \mathbf{0}$).
\end{itemize}

\noindent
Crucially, both conditions use the \textbf{same exploration rate}.
This isolates the effect of prior initialisation: the only
difference is whether the bandit's arm estimates are informed
by offline data or start from uninformative defaults.

\paragraph{Protocol.}
The router is deployed \textbf{from cold start} directly on the
held-out test split ($n{=}1{,}824$) at $\lambda_c{=}0.15$.
There is no pre-training phase---the router sees its first prompt
at $t{=}1$ and learns purely from online partial feedback.
Results are averaged over 20~independent seeds with matched
shuffle order between conditions.
The validation split ($n{=}1{,}785$) is reserved for hyperparameter
selection and the training split ($n{=}8{,}374$) is used only for
offline prior construction.

\paragraph{Cold-start regret (Figure~\ref{fig:warmup_ablation_k2}
and Table~\ref{tab:regret_reduction}).}
Three findings emerge from the cold-start evaluation:

\begin{enumerate}[leftmargin=*,noitemsep]
  \item \textbf{Immediate routing quality.}
        BanditGPT routes effectively from prompt~1, starting
        with windowed reward near~$0.92$ (panel~b), while Tabula
        Rasa begins below~$0.86$ and exhibits wide seed-level
        variance.
        In the first 100~prompts---the critical onboarding
        window---warmup priors reduce cumulative regret by
        \textbf{61.3\%} ($3.0 \pm 0.2$ vs.\ $7.7 \pm 0.2$).
  \item \textbf{Persistent regret gap.}
        Over the full 1{,}824-prompt horizon, warmup priors
        reduce cumulative regret by \textbf{19.6\%}
        ($69.3 \pm 0.5$ vs.\ $86.2 \pm 1.5$; panel~a).
        The gap opens from prompt~1 and widens monotonically:
        priors front-load the benefit but never lose it,
        because cumulative regret is irreversible.
  \item \textbf{Faster convergence.}
        BanditGPT reaches steady-state routing quality in
        $259 \pm 83$~prompts, compared to $487 \pm 90$ for
        Tabula Rasa---a 47\% reduction in the onboarding period.
\end{enumerate}

\paragraph{Interpretation.}
The cold-start protocol reveals that warmup priors function
primarily as a \textbf{deployment safety mechanism} rather than
a steady-state performance booster.
After sufficient online experience (as in Figure~\ref{fig:pareto_k2}),
both conditions converge to comparable routing quality.
The prior's value is concentrated in the first few hundred prompts,
where it eliminates the costly exploration phase that an uninformed
bandit must endure.

This finding directly motivates the system design in
Figure~\ref{fig:adaptive_drift}: when the drift detector triggers a
reset, the router re-enters a cold-start state.
Warmup priors ensure that the post-reset onboarding period
is brief and well-calibrated, avoiding the ${\sim}500$-prompt
exploration cost that a tabula rasa reset would incur.

\paragraph{Information budget.}
The warmup priors are constructed from \emph{full-information}
rewards---observing all arms' quality on every training prompt---while
Tabula Rasa learns only from partial (bandit) feedback.
BanditGPT therefore has a strictly larger information budget.
We note two points in defence of this comparison:

\begin{enumerate}[leftmargin=*,noitemsep]
  \item \textbf{Standard evaluation assumption.}
        The same full-information data is available to all supervised
        baselines in Figure~\ref{fig:pareto_k2} and is the standard
        setting in LLM routing benchmarks~\cite{ong2024routellm,
        panda2025pilot}.
  \item \textbf{Realistic production analogue.}
        In deployment, the full-information phase corresponds to an
        initial A/B testing or uniform exploration period where all
        candidate models are evaluated.  This is standard practice
        before engaging a contextual
        bandit~\cite{kveton2023warmstart} and constitutes a bounded,
        one-time data-collection cost.
\end{enumerate}

\paragraph{Relationship to prior work.}
\citet{kveton2023warmstart} prove that warm-started Linear Thompson
Sampling achieves tighter regret bounds than cold-start under
standard assumptions.
\citet{pan2024cbli} show that even approximate priors reduce
online regret in contextual bandits.
Our results complement these findings with a key practical insight:
the warm-start benefit concentrates in the \emph{onboarding period}
and manifests as a deployment safety property---eliminating the
${\sim}500$-prompt exploration cost---rather than improving
asymptotic routing quality.
This framing is actionable: operators should invest in offline
prior construction when cold-start cost is operationally significant
(e.g., customer-facing systems, high per-query API costs), and
can skip it when an extended burn-in period is acceptable.
